\chapter{单自由度系统强迫振动——一般激励}
\intro{本章将研究任意时间历程激励下单自由度系统的响应计算方法。我们从周期激励的 Fourier 级数展开出发，引入单位脉冲响应函数和卷积积分（Duhamel 积分），建立时域求解框架;随后讨论频域方法（Fourier 变换）和冲击响应谱，最后介绍数值求解方法。这些内容构成了振动分析和结构动力学的核心计算工具。}

\section{周期激励的 Fourier 级数解法}

\subsection{Fourier 级数的基本概念}

在实际工程中，许多激励力具有周期性，如旋转机械的不平衡力、往复式发动机的惯性力等。任何满足 Dirichlet 条件的周期函数 $F(t)$（周期 $T = 2\pi/\omega$）都可以展开为 Fourier 级数：

\begin{myformula}
\[
F(t) = \frac{a_0}{2} + \sum_{n=1}^{\infty} \bigl( a_n \cos n\omega t + b_n \sin n\omega t \bigr)
\]
\end{myformula}

其中 Fourier 系数由以下积分给出：

\begin{myformula}
\[
a_n = \frac{2}{T}\int_{0}^{T} F(t) \cos n\omega t \,dt,\quad n=0,1,2,\ldots
\]
\[
b_n = \frac{2}{T}\int_{0}^{T} F(t) \sin n\omega t \,dt,\quad n=1,2,3,\ldots
\]
\end{myformula}

式中 $\omega = 2\pi/T$ 为基频，$n\omega$ 为第 $n$ 次谐波频率。$a_0/2$ 为激励的静态分量（平均值）。

工程中对周期性激励宜采用 Fourier 级数展开的原因在于：任意复杂的周期激励可分解为无穷多个简谐力分量之和，这就将任意周期激励的响应问题转化为多个简谐激励响应问题的叠加——这正是线性系统的核心优势。

\subsection{叠加原理与总响应}

考虑受迫振动方程：

\begin{myformula}
\[
m\ddot{x} + c\dot{x} + kx = F(t) = \frac{a_0}{2} + \sum_{n=1}^{\infty} \bigl( a_n \cos n\omega t + b_n \sin n\omega t \bigr)
\]
\end{myformula}

由叠加原理，稳态响应为各谐波分量响应的代数和。对于常数项 $a_0/2$，其稳态响应为 $x_0 = a_0/(2k)$。对于第 $n$ 次谐波分量 $a_n \cos n\omega t + b_n \sin n\omega t$，可将其写为单一正弦（或余弦）形式：

\[
a_n \cos n\omega t + b_n \sin n\omega t = F_n \sin(n\omega t + \phi_{Fn})
\]

其中 $F_n = \sqrt{a_n^2 + b_n^2}$，$\phi_{Fn} = \arctan(a_n/b_n)$（需注意象限）。

该谐波分量作用下稳态响应的幅值和相位为：

\begin{myformula}
\[
X_n = \frac{F_n/k}{\sqrt{(1 - r_n^2)^2 + (2\zeta r_n)^2}},\quad r_n = \frac{n\omega}{\omega_n}
\]
\[
\phi_n = \arctan\frac{2\zeta r_n}{1 - r_n^2}
\]
\end{myformula}

因此，总稳态响应为：

\begin{myformula}
\[
x(t) = \frac{a_0}{2k} + \sum_{n=1}^{\infty} \frac{\sqrt{a_n^2 + b_n^2}/k}{\sqrt{(1 - r_n^2)^2 + (2\zeta r_n)^2}} \sin\bigl( n\omega t + \phi_{Fn} - \phi_n \bigr)
\]
\end{myformula}

\subsection{截断与收敛性讨论}

实际计算中只能取有限项近似。Fourier 级数的收敛速度取决于 $F(t)$ 的光滑性：

\begin{itemize}
\item 若 $F(t)$ 连续且分段光滑，则 Fourier 系数 $a_n, b_n \sim O(1/n^2)$，级数收敛较快。
\item 若 $F(t)$ 存在跳跃间断（如方波），则 $a_n, b_n \sim O(1/n)$，收敛缓慢，且出现 Gibbs 现象（间断点附近约 9\% 的过冲）。
\end{itemize}

在振动问题中，由于系统的低通滤波特性（$|H(\omega)| \to 0$ 当 $\omega \gg \omega_n$），高频分量的贡献很小，一般截取前 $N$ 阶即可满足工程精度。通常取 $N$ 使得 $r_N = N\omega/\omega_n \geq 3\sim5$ 即可。

\subsection{周期激励响应的 TikZ 示意图}

\begin{tikzpicture}[scale=0.9]
  \def\xmax{10}
  \def\ymax{3}
  % 周期方波激励及其 Fourier 近似
  \draw[-Stealth] (0,0) -- (\xmax,0) node[below] {$t$};
  \draw[-Stealth] (0,-\ymax) -- (0,\ymax) node[left] {$F(t)$};
  % 方波
  \draw[thick,blue] (0,1.5) -- (1.5,1.5) -- (1.5,-1.5) -- (3.0,-1.5) -- (3.0,1.5) -- (4.5,1.5) -- (4.5,-1.5) -- (6.0,-1.5) -- (6.0,1.5) -- (7.5,1.5) -- (7.5,-1.5) -- (9.0,-1.5);
  \node[blue,above] at (2.25,1.8) {方波激励};
  % 1项谐波
  \draw[red,dashed,domain=0:9,samples=100] plot (\x,{1.9099*sin(0.698*\x r)});
  \node[red,right] at (9.2,0.5) {$n=1$};
  % 系统响应示意
  \draw[thick,green!60!black,domain=0:9,samples=200] plot (\x,{0.8*sin(0.698*\x r - 30)});
  \node[green!60!black,right] at (9.2,-1.5) {稳态响应};
\end{tikzpicture}

\section{单位脉冲响应}

\subsection{Dirac $\delta$ 函数的定义与性质}

Dirac $\delta$ 函数（或单位脉冲函数）是振动分析与信号处理中最基本的"原子"激励。它不是通常意义上的函数，而是一个广义函数（分布），由以下两个条件定义：

\begin{myformula}
\[
\delta(t - \tau) = 0,\quad t \neq \tau;\qquad \int_{-\infty}^{\infty} \delta(t - \tau)\,dt = 1
\]
\end{myformula}

其核心性质是筛选性质：

\begin{myformula}
\[
\int_{-\infty}^{\infty} f(t)\,\delta(t - \tau)\,dt = f(\tau)
\]
\end{myformula}

物理上，$\delta(t-\tau)$ 表示在 $t=\tau$ 时刻作用的、幅值无穷大但冲量（积分值）等于 1 的理想化力。它的量纲为 $[\text{时间}]^{-1}$。

$\delta$ 函数可视为一系列普通函数序列的极限。例如矩形脉冲序列：

\[
\delta(t) = \lim_{\varepsilon \to 0} \frac{1}{\varepsilon}\bigl[ H(t) - H(t-\varepsilon) \bigr]
\]

或高斯函数的极限：

\[
\delta(t) = \lim_{\sigma \to 0} \frac{1}{\sqrt{2\pi}\,\sigma} e^{-t^2/(2\sigma^2)}
\]

$\delta$ 函数的导数 $\delta'(t)$ 也具有物理意义，表示力偶（一对大小相等方向相反的脉冲），在弯矩激励问题中有应用。

\subsection{脉冲响应函数 $h(t)$ 的推导}

考虑在 $t=0$ 时刻受到单位脉冲作用的欠阻尼单自由度系统：

\[
m\ddot{x} + c\dot{x} + kx = \delta(t)
\]

在 $t=0$ 处积分：

\[
\int_{0^-}^{0^+} (m\ddot{x} + c\dot{x} + kx)\,dt = \int_{0^-}^{0^+} \delta(t)\,dt = 1
\]

由于 $x$ 在 $t=0$ 处连续（位移不能突变），而速度可以突变。在无穷小时间区间内：
\begin{itemize}
\item $\int_{0^-}^{0^+} kx\,dt = 0$（因为 $x$ 连续且区间长度 $\to 0$）
\item $\int_{0^-}^{0^+} c\dot{x}\,dt \approx c\int_{0^-}^{0^+} \dot{x}\,dt = c \cdot \Delta x \to 0$（位移增量有限，但在 $\varepsilon \to 0$ 时 $\Delta x \to 0$）
\item $\int_{0^-}^{0^+} m\ddot{x}\,dt = m[\dot{x}(0^+) - \dot{x}(0^-)] = 1$
\end{itemize}

因此初始条件为：

\[
x(0^+) = 0,\quad \dot{x}(0^+) = \frac{1}{m}
\]

这相当于一个具有初速度 $1/m$ 的自由振动问题，其解就是单位脉冲响应函数（Impulse Response Function, IRF）：

\begin{myformula}
\[
h(t) = \frac{1}{m\omega_d} e^{-\zeta\omega_n t} \sin\omega_d t,\quad t > 0
\]
\end{myformula}

其中 $\omega_d = \omega_n\sqrt{1-\zeta^2}$ 为有阻尼固有频率。$h(t)=0$ 当 $t<0$（因果性）。

\subsection{脉冲响应函数的性质与物理意义}

\begin{enumerate}
\item \textbf{因果性}：$h(t)=0$ 对 $t<0$。系统在激励作用之前不能有响应。

\item \textbf{线性时不变性}：若脉冲作用时刻为 $t=\tau$，则响应为 $h(t-\tau)$。

\item \textbf{系统的"记忆"}：$h(t)$ 以指数衰减的正弦函数形式描述了系统如何"记住"并逐渐"遗忘"一次脉冲扰动。衰减速率由 $\zeta\omega_n$ 决定，振荡频率为 $\omega_d$。

\item \textbf{脉冲响应面积}：稳态增益 $\int_0^\infty h(t)\,dt = 1/k$，即单位脉冲的长期位移响应等于静柔度。

\item \textbf{传递函数的关系}：$h(t)$ 的 Laplace 变换即为传递函数 $H(s) = 1/(ms^2+cs+k)$。
\end{enumerate}

\begin{tikzpicture}[scale=0.9]
  \def\xmax{8}
  % 脉冲响应函数曲线
  \draw[-Stealth] (0,-2.5) -- (0,2.8) node[above] {$h(t)$};
  \draw[-Stealth] (0,0) -- (\xmax,0) node[below] {$t$};
  % 画 h(t) 曲线，欠阻尼情况
  \draw[thick,blue,domain=0:7.5,samples=300] plot (\x,{2.0*exp(-0.3*\x)*sin(2.0*\x r)});
  \draw[thick,red,dashed,domain=0:7.5,samples=100] plot (\x,{2.0*exp(-0.3*\x)});
  \draw[thick,red,dashed,domain=0:7.5,samples=100] plot (\x,{-2.0*exp(-0.3*\x)});
  \node[blue,right] at (7.8,0.6) {$h(t)$};
  \node[red,right] at (7.8,2.2) {$\pm e^{-\zeta\omega_n t}/(m\omega_d)$};
  % 标记特征量
  \draw[thin,dotted] (0.79,2.0) -- (0.79,0) node[below] {$T_d$};
  \node at (4,-2.8) {$T_d = \frac{2\pi}{\omega_d}$：有阻尼周期};
\end{tikzpicture}

\section{任意激励的响应——卷积积分（Duhamel 积分）}

\subsection{基本思想：脉冲序列的叠加}

任意力 $F(t)$ 可以看作一系列时间上连续的、不同幅值脉冲的叠加。在时刻 $t=\tau$、宽度为 $d\tau$ 的脉冲，其冲量为 $F(\tau)\,d\tau$。根据线性时不变系统的性质，该微元脉冲在 $t$ 时刻（$t>\tau$）产生的响应微元为：

\[
dx(t) = F(\tau)\,d\tau \cdot h(t-\tau)
\]

将所有 $\tau \in [0, t]$ 上的脉冲微元贡献累加，得到：

\begin{myformula}
\[
x(t) = \int_{0}^{t} F(\tau)\,h(t-\tau)\,d\tau
\]
\end{myformula}

这就是 \textbf{Duhamel 积分}（或卷积积分）。它是时域中计算线性系统任意激励响应的基本公式。其交换对称形式为：

\begin{myformula}
\[
x(t) = \int_{0}^{t} F(t-\tau)\,h(\tau)\,d\tau
\]
\end{myformula}

两者等价，在数值计算时可选择便于离散化的形式。

将 $h(t)$ 的具体表达式代入，得到欠阻尼系统的 Duhamel 积分完整形式：

\begin{myformula}
\[
x(t) = \frac{1}{m\omega_d} \int_{0}^{t} F(\tau)\, e^{-\zeta\omega_n (t-\tau)} \sin\bigl[ \omega_d (t-\tau) \bigr]\,d\tau
\]
\end{myformula}

该式给出了零初始条件下任意力激励的响应。若存在非零初始条件，需叠加自由振动项：

\[
x(t) = e^{-\zeta\omega_n t}\Bigl( x_0 \cos\omega_d t + \frac{\dot{x}_0 + \zeta\omega_n x_0}{\omega_d}\sin\omega_d t \Bigr) + \frac{1}{m\omega_d}\int_{0}^{t} F(\tau) e^{-\zeta\omega_n(t-\tau)} \sin\omega_d(t-\tau)\,d\tau
\]

\subsection{卷积积分的物理图示}

\begin{tikzpicture}[scale=0.75]
  % 上子图：激励力 F(t)
  \draw[-Stealth] (0,0) -- (12,0) node[below] {$t$};
  \draw[-Stealth] (0,0) -- (0,3.5) node[above] {$F(\tau)$};
  \draw[thick,blue] (0,0) -- (1.5,2.5) -- (3,1) -- (4.5,3) -- (6,0.5) -- (7.5,2) -- (9,1.5) -- (10.5,2.8) -- (12,1);
  % 标记脉冲微元
  \draw[red,thick] (4.5,0) -- (4.5,3);
  \draw[red,thick] (4.7,3) -- (4.7,3.5);
  \node[red,above] at (4.5,3.6) {$F(\tau)d\tau$};
  \draw[->,red] (4.6,3.3) -- (5.5,3.3) node[right] {冲量微元};
  \node at (6,-0.8) {激励力 $F(\tau)$ 分解为脉冲序列};

  % 下子图：响应 x(t)
  \begin{scope}[yshift=-5cm]
    \draw[-Stealth] (0,0) -- (12,0) node[below] {$t$};
    \draw[-Stealth] (0,-2.5) -- (0,2.8) node[above] {$x(t)$};
    % 画各个脉冲元产生的响应
    \draw[domain=0.5:11.5,samples=200,blue!60] plot (\x,{0.5*exp(-0.4*(\x-0.5))*sin(3*(\x-0.5) r)});
    \draw[domain=2:11.5,samples=200,green!60!black] plot (\x,{0.8*exp(-0.4*(\x-2.0))*sin(3*(\x-2.0) r)});
    \draw[domain=4.5:11.5,samples=200,red!80] plot (\x,{1.2*exp(-0.4*(\x-4.5))*sin(3*(\x-4.5) r)});
    \draw[domain=7:11.5,samples=200,orange] plot (\x,{0.6*exp(-0.4*(\x-7.0))*sin(3*(\x-7.0) r)});
    % 总和
    \draw[domain=0:11.5,samples=400,thick,red] plot (\x,{
      0.5*exp(-0.4*max(0,\x-0.5))*sin(3*max(0,\x-0.5) r)+
      0.8*exp(-0.4*max(0,\x-2.0))*sin(3*max(0,\x-2.0) r)+
      1.2*exp(-0.4*max(0,\x-4.5))*sin(3*max(0,\x-4.5) r)+
      0.6*exp(-0.4*max(0,\x-7.0))*sin(3*max(0,\x-7.0) r)
    });
    \node[red,right] at (11.8,1.2) {总响应};
    \node at (6,-3.2) {各脉冲产生的响应叠加 = 卷积积分};
  \end{scope}
\end{tikzpicture}

\subsection{阶跃响应（Step Response）}

阶跃力 $F(t) = F_0 H(t)$（$H(t)$ 为 Heaviside 阶跃函数）是工程中最常见的瞬态激励之一。将其代入 Duhamel 积分：

\[
x(t) = \frac{F_0}{m\omega_d} \int_{0}^{t} e^{-\zeta\omega_n (t-\tau)} \sin\omega_d (t-\tau)\,d\tau
\]

作变量代换 $u = t-\tau$，$du = -d\tau$：

\[
x(t) = \frac{F_0}{m\omega_d} \int_{0}^{t} e^{-\zeta\omega_n u} \sin\omega_d u\,du
\]

利用积分公式 $\int e^{ax}\sin bx\,dx = \frac{e^{ax}}{a^2+b^2}(a\sin bx - b\cos bx)$，其中 $a = -\zeta\omega_n$，$b = \omega_d$，并注意到 $a^2+b^2 = \zeta^2\omega_n^2 + \omega_n^2(1-\zeta^2) = \omega_n^2$，得到：

\begin{myformula}
\[
x(t) = \frac{F_0}{k} \Bigl[ 1 - e^{-\zeta\omega_n t}\Bigl( \cos\omega_d t + \frac{\zeta}{\sqrt{1-\zeta^2}} \sin\omega_d t \Bigr) \Bigr]
\]
\end{myformula}

\textbf{物理意义}：
\begin{itemize}
\item 稳态值 $x_{\text{ss}} = F_0/k$，等于静力作用下的位移（柔度 $\times$ 力）。
\item 瞬态项在 $t \to \infty$ 时衰减为零，衰减速率由 $\zeta\omega_n$ 决定。
\item 最大超调量（过冲）与 $\zeta$ 密切相关。当 $\zeta=0$ 时，$x_{\max} = 2F_0/k$（100\% 超调）;随 $\zeta$ 增大，超调减小。
\item 上升时间、峰值时间、稳定时间等时域性能指标可直接由该式导出。
\end{itemize}

\subsection{斜坡响应（Ramp Response）}

斜坡力 $F(t) = \alpha t H(t)$（$\alpha$ 为加载速率）对应缓慢加载或匀速加载工况。代入 Duhamel 积分并积分，得：

\begin{myformula}
\[
x(t) = \frac{\alpha}{k} \Bigl[ t - \frac{2\zeta}{\omega_n} + e^{-\zeta\omega_n t}\Bigl( \frac{2\zeta}{\omega_n}\cos\omega_d t + \frac{2\zeta^2-1}{\omega_d}\sin\omega_d t \Bigr) \Bigr]
\]
\end{myformula}

其稳态响应为 $x_{\text{ss}}(t) = \frac{\alpha}{k}\bigl(t - \frac{2\zeta}{\omega_n}\bigr)$，即与输入斜率相同的斜坡，但有一个固定的滞后时间 $2\zeta/\omega_n$。这一滞后反映了阻尼的耗能特性——系统需要"追上"不断增长的激励。

\subsection{Duhamel 积分数值计算的 TikZ 示意图}

\begin{tikzpicture}[scale=0.9]
  \def\xmax{8}
  % 离散卷积示意
  \draw[-Stealth] (0,0) -- (\xmax,0) node[below] {$\tau$};
  \draw[-Stealth] (0,0) -- (0,3.5) node[above] {};
  % F(tau) 采样
  \foreach \i in {0,1,...,7} {
    \pgfmathsetmacro{\x}{0.5+1.0*\i}
    \pgfmathsetmacro{\y}{1.5+0.3*rand}
    \fill[blue] (\x,\y) circle (2pt);
    \draw[blue,dashed] (\x,0) -- (\x,\y);
  }
  \node[blue] at (4.5,3.2) {$F(\tau_k)$ 采样点};
  % 数值积分示意
  \draw[->,thick] (8.5,1.5) -- (9.5,1.5) node[right] {$\displaystyle\sum_k F(\tau_k)h(t-\tau_k)\Delta\tau$};
  \node at (4.5,-0.8) {数值 Duhamel 积分：梯形法则离散化};
\end{tikzpicture}

\section{频域方法——Fourier 变换}

\subsection{从时域到频域}

在周期激励中我们用 Fourier 级数将周期函数分解为离散频谱的谐波之和。对于非周期激励——即绝大多数实际瞬态激励（冲击、地震波、风荷载等）——需要使用 Fourier 变换（积分）将其表示为连续频谱。

信号 $F(t)$ 的 Fourier 变换 $\hat{F}(\omega)$ 定义为：

\begin{myformula}
\[
\hat{F}(\omega) = \mathcal{F}\{F(t)\} = \int_{-\infty}^{\infty} F(t)\, e^{-i\omega t}\,dt
\]
\end{myformula}

其逆变换为：

\begin{myformula}
\[
F(t) = \mathcal{F}^{-1}\{\hat{F}(\omega)\} = \frac{1}{2\pi} \int_{-\infty}^{\infty} \hat{F}(\omega)\, e^{i\omega t}\,d\omega
\]
\end{myformula}

$\hat{F}(\omega)$ 是复值函数，其模 $|\hat{F}(\omega)|$ 为幅值谱，辐角 $\angle\hat{F}(\omega)$ 为相位谱。$|\hat{F}(\omega)|^2$ 为能量谱密度（Parseval 定理：时域能量 = 频域能量）。

\subsection{频响函数 $H(\omega)$ 与频域求解}

对运动方程 $m\ddot{x} + c\dot{x} + kx = F(t)$ 两边作 Fourier 变换，利用微分性质 $\mathcal{F}\{\ddot{x}\} = (i\omega)^2\hat{X}(\omega) = -\omega^2\hat{X}(\omega)$，$\mathcal{F}\{\dot{x}\} = i\omega\hat{X}(\omega)$，得到：

\[
(-m\omega^2 + ic\omega + k) \hat{X}(\omega) = \hat{F}(\omega)
\]

定义 \textbf{频响函数}（Frequency Response Function, FRF）：

\begin{myformula}
\[
H(\omega) = \frac{1}{-m\omega^2 + ic\omega + k} = \frac{1/k}{1 - r^2 + 2i\zeta r},\quad r = \frac{\omega}{\omega_n}
\]
\end{myformula}

则频域中的输入输出关系为：

\begin{myformula}
\[
\hat{X}(\omega) = H(\omega) \cdot \hat{F}(\omega)
\]
\end{myformula}

这是一个极其简洁的关系：\textbf{时域卷积 = 频域乘积}。响应的时域形式可由 Fourier 逆变换求得：

\[
x(t) = \frac{1}{2\pi} \int_{-\infty}^{\infty} H(\omega)\,\hat{F}(\omega)\, e^{i\omega t}\,d\omega
\]

\subsection{$H(\omega)$ 的物理意义}

频响函数 $H(\omega)$ 完整刻画了线性系统的动力学特性：

\begin{enumerate}
\item \textbf{幅频特性}：$|H(\omega)| = \frac{1/k}{\sqrt{(1-r^2)^2 + (2\zeta r)^2}}$，表示输出幅值与输入幅值之比。
\item \textbf{相频特性}：$\angle H(\omega) = -\arctan\frac{2\zeta r}{1-r^2}$，表示输出滞后输入的相位角。
\item \textbf{共振}：当 $r=1$（$\omega = \omega_n$）时，$|H(\omega_n)| = 1/(2k\zeta)$，在小阻尼下幅值极大。
\item \textbf{带宽}：$|H(\omega)|$ 的半功率点（$|H|_{\max}/\sqrt{2}$）之间的频率区间，与阻尼比有关：$\Delta\omega \approx 2\zeta\omega_n$。
\item \textbf{与传递函数的关系}：$H(\omega) = H(s)\big|_{s=i\omega}$，即 FRF 是传递函数沿虚轴的取值。
\end{enumerate}

\subsection{Fourier 变换法与传统 Duhamel 积分的关系}

Fourier 变换法和 Duhamel 积分是同一物理过程在频域和时域的两种等价表示。两者的联系由卷积定理保证。在数值计算中：
\begin{itemize}
\item 若激励有解析的 Fourier 变换且系统简单，频域法更简洁。
\item 若激励为实测时程数据（如地震加速度记录），FFT 算法使频域法极为高效。
\item 若系统非线性或需考虑时变参数，则需回到时域的逐步积分法。
\end{itemize}

\section{冲击响应谱（Shock Response Spectrum）}

\subsection{定义与基本概念}

\textbf{冲击响应谱}（Shock Response Spectrum, SRS）是衡量冲击对结构影响的重要工程工具。其定义为：一系列不同固有频率的单自由度振子在同一冲击激励下，各自的最大响应（位移、速度或加速度）与固有频率的关系曲线。

数学定义：

\begin{myformula}
\[
S_d(\omega_n, \zeta) = \max_{t} |x(t; \omega_n, \zeta)|
\]
\end{myformula}

其中 $x(t; \omega_n, \zeta)$ 是固有频率为 $\omega_n$、阻尼比为 $\zeta$ 的单自由度系统在指定冲击输入下的位移响应。

常用的三种冲击响应谱：
\begin{itemize}
\item \textbf{位移谱} $S_d$：最大相对位移
\item \textbf{伪速度谱} $S_v = \omega_n S_d$：与最大应变能相关
\item \textbf{伪加速度谱} $S_a = \omega_n^2 S_d$：与最大惯性力（基底剪力）相关
\end{itemize}

在中低频段，$S_v$ 近似为常数（等速度区）;在低频段，$S_d$ 趋于冲击的最大位移（等位移区）;在高频段，$S_a$ 趋于冲击的最大加速度（等加速度区）。

\subsection{半正弦脉冲的冲击响应谱}

半正弦脉冲力 $F(t) = F_0\sin(\pi t/t_d)$，$0 \leq t \leq t_d$。系统的最大响应发生在脉冲作用期间或之后，取决于 $t_d/T_n$ 的比值（$T_n = 2\pi/\omega_n$ 为系统固有周期）。

将半正弦脉冲代入 Duhamel 积分，对无阻尼情况可得到闭合形式的响应表达式（分 $t \leq t_d$ 和 $t > t_d$ 两段）。最大响应（位移谱）为：

\begin{myformula}
\[
\frac{S_d}{F_0/k} = 
\begin{cases}
\frac{1}{1 - (2\tau)^2}\bigl[ \sin(\pi t/t_d) - 2\tau\sin(2\pi t/t_d)/(2\tau) \bigr], & t \leq t_d \\[6pt]
\frac{2\tau}{1 - (2\tau)^2} \bigl[ 1 + \cos(2\pi\tau) \bigr]^{1/2}, & t > t_d,\; \tau = t_d/T_n
\end{cases}
\]
\end{myformula}

关键特征：
\begin{itemize}
\item 当 $t_d/T_n \ll 1$（短脉冲），最大响应 $S_d \approx 2I/(m\omega_n)$，其中 $I = \int F(t)dt$ 为冲量。
\item 当 $t_d/T_n \gg 1$（长脉冲），最大响应趋近静力解的两倍 $S_d \approx 2F_0/k$。
\item 最大动态放大因子约为 1.75（出现在 $t_d/T_n \approx 0.7$ 附近）。
\end{itemize}

\subsection{矩形脉冲与三角脉冲的冲击响应谱}

\textbf{矩形脉冲} $F(t) = F_0$，$0 \leq t \leq t_d$。无阻尼下的最大响应：

\[
\frac{S_d}{F_0/k} = 
\begin{cases}
2\sin(\pi t_d/T_n), & t_d/T_n \leq 1/2 \\
2, & t_d/T_n > 1/2
\end{cases}
\]

\textbf{三角脉冲}（等腰）$F(t) = F_0(1 - 2|t - t_d/2|/t_d)$，$0 \leq t \leq t_d$。其冲击响应谱的峰值低于矩形脉冲，因为其"圆滑"的力历程减小了高次谐波含量。

三种脉冲响应谱的比较表明：
\begin{itemize}
\item 矩形脉冲（最尖锐的加载/卸载）产生最大的动态放大因子（约 2.0）。
\item 半正弦脉冲的动态放大因子居中（约 1.75）。
\item 三角脉冲（最平滑）的动态放大因子最小（约 1.5）。
\item 这反映了一个一般性原理：冲击的尖锐程度（高频含量）决定最大响应。
\end{itemize}

\subsection{冲击响应谱的 TikZ 图}

\begin{tikzpicture}[scale=0.85]
  % 三脉冲的 SRS 对比
  \draw[-Stealth] (0,0) -- (12,0) node[below] {$t_d / T_n$};
  \draw[-Stealth] (0,0) -- (0,5) node[above] {$S_d / (F_0/k)$};
  % 矩形脉冲 SRS
  \draw[thick,blue,domain=0.01:5,samples=200] plot (\x*2.2,{min(2*sin(3.14159*\x r), 2.0)*1.2});
  \node[blue,right] at (11,3.5) {矩形脉冲};
  % 半正弦脉冲 SRS
  \draw[thick,red,dashed,domain=0.01:5,samples=200] plot (\x*2.2,{min(1.75*sin(2.5*\x r), 1.75)*1.2});
  \node[red,right] at (11,2.5) {半正弦脉冲};
  % 三角脉冲 SRS
  \draw[thick,green!60!black,dotted,domain=0.01:5,samples=200] plot (\x*2.2,{min(1.5*sin(2.0*\x r), 1.5)*1.2});
  \node[green!60!black,right] at (11,1.5) {三角脉冲};
  % 标注
  \draw[dashed,gray] (0,2.4) -- (11,2.4) node[right,gray] {$2.0$};
  \draw[dashed,gray] (0,2.1) -- (11,2.1);
  \draw[dashed,gray] (0,1.8) -- (11,1.8);
\end{tikzpicture}

\section{数值方法}

\subsection{直接数值积分法的基本思路}

当激励 $F(t)$ 为实测时间序列或过于复杂无法解析积分时，需要数值方法求解运动方程。数值方法的基本思路是将时间离散为步长 $\Delta t$ 的时间点序列 $t_0, t_1, \ldots, t_N$，在每个时间步内对运动方程进行近似积分，递推求解响应。

运动方程在 $t_{n+1}$ 时刻成立：
\[
m\ddot{x}_{n+1} + c\dot{x}_{n+1} + k x_{n+1} = F_{n+1}
\]

需要补充位移、速度、加速度之间的关系式（差分格式）才能递推求解。

\subsection{Newmark-$\beta$ 法}

Newmark-$\beta$ 法是最常用的直接积分法之一，其基本假设为速度和位移的递推关系：

\begin{myformula}
\[
\dot{x}_{n+1} = \dot{x}_n + (1-\gamma)\Delta t\,\ddot{x}_n + \gamma\Delta t\,\ddot{x}_{n+1}
\]
\[
x_{n+1} = x_n + \Delta t\,\dot{x}_n + \Bigl(\frac{1}{2} - \beta\Bigr)\Delta t^2\,\ddot{x}_n + \beta\Delta t^2\,\ddot{x}_{n+1}
\]
\end{myformula}

参数 $\beta$ 和 $\gamma$ 控制方法的精度和稳定性：

\begin{itemize}
\item $\gamma = 1/2$ 保证二阶精度。
\item $\beta = 1/4$（平均加速度法/梯形法则）：无条件稳定，最常用。
\item $\beta = 1/6$（线性加速度法）：有条件稳定，$\Delta t/T_n \leq 0.551$。
\item $\beta = 1/12$（Fox-Goodwin 法）：有条件稳定。
\end{itemize}

算法步骤（$\beta=1/4$，$\gamma=1/2$）：

\begin{enumerate}
\item 初始计算：由初始条件 $x_0$, $\dot{x}_0$ 和运动方程求 $\ddot{x}_0 = (F_0 - c\dot{x}_0 - kx_0)/m$。
\item 形成有效刚度矩阵：$\hat{k} = k + \frac{\gamma}{\beta\Delta t}c + \frac{1}{\beta\Delta t^2}m$。
\item 对每个时间步 $n = 0, 1, \ldots, N-1$：
  \begin{itemize}
  \item 计算有效载荷：$\hat{F}_{n+1} = F_{n+1} + m\bigl(\frac{1}{\beta\Delta t^2}x_n + \frac{1}{\beta\Delta t}\dot{x}_n + (\frac{1}{2\beta}-1)\ddot{x}_n\bigr) + c\bigl(\frac{\gamma}{\beta\Delta t}x_n + (\frac{\gamma}{\beta}-1)\dot{x}_n + \frac{\Delta t}{2}(\frac{\gamma}{\beta}-2)\ddot{x}_n\bigr)$。
  \item 求解位移：$x_{n+1} = \hat{F}_{n+1}/\hat{k}$。
  \item 更新加速度：$\ddot{x}_{n+1} = \frac{1}{\beta\Delta t^2}(x_{n+1} - x_n) - \frac{1}{\beta\Delta t}\dot{x}_n - (\frac{1}{2\beta}-1)\ddot{x}_n$。
  \item 更新速度：$\dot{x}_{n+1} = \dot{x}_n + (1-\gamma)\Delta t\ddot{x}_n + \gamma\Delta t\ddot{x}_{n+1}$。
  \end{itemize}
\end{enumerate}

\subsection{Runge-Kutta 法}

将二阶运动方程化为一阶状态空间方程组：

\[
\frac{d}{dt} \begin{pmatrix} x \\ \dot{x} \end{pmatrix} = 
\begin{pmatrix} \dot{x} \\ \frac{1}{m}(F(t) - c\dot{x} - kx) \end{pmatrix}
\]

即 $\dot{\mathbf{y}} = \mathbf{f}(t, \mathbf{y})$，其中 $\mathbf{y} = [x, \dot{x}]^{\top}$。对该一阶系统可直接应用标准的四阶 Runge-Kutta 法（RK4）：

\begin{myformula}
\[
\mathbf{k}_1 = \Delta t \cdot \mathbf{f}(t_n, \mathbf{y}_n)
\]
\[
\mathbf{k}_2 = \Delta t \cdot \mathbf{f}(t_n + \tfrac{\Delta t}{2}, \mathbf{y}_n + \tfrac{1}{2}\mathbf{k}_1)
\]
\[
\mathbf{k}_3 = \Delta t \cdot \mathbf{f}(t_n + \tfrac{\Delta t}{2}, \mathbf{y}_n + \tfrac{1}{2}\mathbf{k}_2)
\]
\[
\mathbf{k}_4 = \Delta t \cdot \mathbf{f}(t_n + \Delta t, \mathbf{y}_n + \mathbf{k}_3)
\]
\[
\mathbf{y}_{n+1} = \mathbf{y}_n + \frac{1}{6}(\mathbf{k}_1 + 2\mathbf{k}_2 + 2\mathbf{k}_3 + \mathbf{k}_4)
\]
\end{myformula}

RK4 具有四阶全局精度（$O(\Delta t^4)$），是显式方法中最常用的。其优点在于不需迭代求解，编程实现简便;缺点是对于刚性系统（高阻尼或高频），稳定步长受限。

\subsection{数值方法的 TikZ 示意图}

\begin{tikzpicture}[scale=0.9]
  \def\xmax{8}
  \def\ymax{3.5}
  % 数值积分步进示意
  \draw[-Stealth] (0,0) -- (\xmax,0) node[below] {$t$};
  \draw[-Stealth] (0,-2) -- (0,\ymax) node[above] {$x(t)$};
  % 精确解曲线
  \draw[thick,blue,domain=0:7.5,samples=300] plot (\x,{exp(-0.2*\x)*sin(2.0*\x r)+1.5});
  \node[blue,right] at (7.8,2.5) {精确解};
  % 数值解采样点
  \foreach \i in {0,1,...,14} {
    \pgfmathsetmacro{\t}{0.5*\i}
    \pgfmathsetmacro{\y}{exp(-0.2*\t)*sin(40*\t)+1.5+0.05*rand}
    \fill[red] (\t,\y) circle (2pt);
  }
  % 步长标注
  \draw[<->] (1.0,-0.3) -- (1.5,-0.3) node[midway,below] {$\Delta t$};
  \node[red,right] at (7.8,1.0) {数值解 $t_n$};
\end{tikzpicture}

\section{Python 代码实现}

\subsection{Duhamel 积分的数值计算}

\begin{mycode}{Duhamel积分数值计算}
\begin{lstlisting}[language=Python]
import numpy as np
import matplotlib.pyplot as plt

def duhamel_response(F, t, m, k, c, x0=0.0, v0=0.0):
    """
    用梯形法则数值计算 Duhamel 积分响应
    F: 激励力数组 (N)
    t: 时间数组 (s)
    m, k, c: 质量, 刚度, 阻尼
    x0, v0: 初始位移和速度
    """
    omega_n = np.sqrt(k / m)       # 固有频率
    zeta = c / (2 * np.sqrt(m * k)) # 阻尼比
    omega_d = omega_n * np.sqrt(1 - zeta**2) if zeta < 1 else 0.0

    dt = t[1] - t[0]
    n = len(t)
    x = np.zeros(n)
    v = np.zeros(n)
    x[0], v[0] = x0, v0

    # 脉冲响应函数 (欠阻尼)
    def h(tau):
        if tau <= 0:
            return 0.0
        return np.exp(-zeta * omega_n * tau) * np.sin(omega_d * tau) / (m * omega_d)

    for i in range(1, n):
        # 梯形法则离散化
        integ_sum = 0.0
        for j in range(i + 1):
            tau = t[i] - t[j]
            fj = F[j]
            if j == 0 or j == i:
                integ_sum += 0.5 * fj * h(tau)
            else:
                integ_sum += fj * h(tau)
        integ_sum *= dt

        # 自由振动项
        edt = np.exp(-zeta * omega_n * t[i])
        cos_term = x0 * np.cos(omega_d * t[i])
        sin_term = (v0 + zeta * omega_n * x0) / omega_d * np.sin(omega_d * t[i])
        x[i] = edt * (cos_term + sin_term) + integ_sum

    return x
\end{lstlisting}
\end{mycode}

\subsection{Newmark-$\beta$ 法求解}

\begin{mycode}{Newmark-$\beta$法求解单自由度系统}
\begin{lstlisting}[language=Python]
def newmark_beta(m, k, c, F, dt, x0=0.0, v0=0.0, beta=0.25, gamma=0.5):
    """
    Newmark-beta 法求解 m*x'' + c*x' + k*x = F(t)
    beta=0.25, gamma=0.5: 平均加速度法 (无条件稳定)
    """
    n = len(F)
    x = np.zeros(n)
    v = np.zeros(n)
    a = np.zeros(n)

    x[0], v[0] = x0, v0
    a[0] = (F[0] - c * v[0] - k * x[0]) / m

    # 有效刚度 (单自由度时为标量)
    k_eff = k + gamma * c / (beta * dt) + m / (beta * dt**2)

    for i in range(n - 1):
        # 预测步
        dp_v = v[i] + (1 - gamma) * dt * a[i]  # 速度预测
        dp_x = x[i] + dt * v[i] + (0.5 - beta) * dt**2 * a[i]  # 位移预测

        # 有效载荷
        F_eff = (F[i+1] +
                 m * (x[i] / (beta * dt**2) + v[i] / (beta * dt) +
                      (0.5 / beta - 1) * a[i]) +
                 c * (dp_x * gamma / (beta * dt) +
                      dp_v * (gamma / beta - 1)))

        # 实际上对于平均加速度法:
        a_coeff = m / (beta * dt**2) + gamma * c / (beta * dt)
        F_pred = (F[i+1] +
                  a_coeff * x[i] +
                  (m / (beta * dt) + c * (gamma / beta - 1)) * v[i] +
                  (m * (0.5 / beta - 1) + c * dt * (0.5 * gamma / beta - 1)) * a[i])

        x[i+1] = F_pred / k_eff
        a[i+1] = (x[i+1] - x[i]) / (beta * dt**2) - v[i] / (beta * dt) - (0.5 / beta - 1) * a[i]
        v[i+1] = v[i] + (1 - gamma) * dt * a[i] + gamma * dt * a[i+1]

    return x, v, a

# 示例：半正弦脉冲激励
dt = 0.01; T = 5.0; t = np.arange(0, T, dt)
td = 0.5; F0 = 1.0
F = np.where((t >= 0) & (t <= td), F0 * np.sin(np.pi * t / td), 0.0)

m, k, c = 1.0, 40.0, 0.4
x_newmark, _, _ = newmark_beta(m, k, c, F, dt)
\end{lstlisting}
\end{mycode}

\subsection{Runge-Kutta 法求解}

\begin{mycode}{Runge-Kutta法求解单自由度系统}
\begin{lstlisting}[language=Python]
def rk4_sdof(m, k, c, F_func, t_span, y0, dt):
    """
    四阶 Runge-Kutta 法求解 m*x'' + c*x' + k*x = F(t)
    F_func: 激励力函数 F(t)
    t_span: (t_start, t_end)
    y0: [x0, v0]
    """
    t_start, t_end = t_span
    n = int((t_end - t_start) / dt) + 1
    t = np.linspace(t_start, t_end, n)

    y = np.zeros((n, 2))  # [x, v]
    y[0] = y0

    for i in range(n - 1):
        ti = t[i]
        yi = y[i]

        def f(t_val, y_val):
            x_i, v_i = y_val
            a_i = (F_func(t_val) - c * v_i - k * x_i) / m
            return np.array([v_i, a_i])

        k1 = dt * f(ti, yi)
        k2 = dt * f(ti + 0.5 * dt, yi + 0.5 * k1)
        k3 = dt * f(ti + 0.5 * dt, yi + 0.5 * k2)
        k4 = dt * f(ti + dt, yi + k3)

        y[i+1] = yi + (k1 + 2*k2 + 2*k3 + k4) / 6.0

    return t, y[:, 0], y[:, 1]
\end{lstlisting}
\end{mycode}

\section{小结}

本章系统建立了单自由度系统在任意激励下的时域和频域求解框架。核心工具链为：Fourier 级数（周期激励）$\to$ 脉冲响应函数 $h(t)$ $\to$ Duhamel 积分（任意激励的时域解）$\to$ Fourier 变换/频响函数（频域解）。冲击响应谱为冲击工程设计提供了简洁的工程工具，而 Newmark-$\beta$ 法和 Runge-Kutta 法等数值方法则为复杂激励下提供了一般的求解途径。这些方法将在多自由度系统（第五章和第六章）中进一步推广。
